\section{为什么需要 Offline RL}

\begin{importantbox}{Offline RL 的定义}
给定一个由未知行为策略 $\pi_\beta$ 收集的\textbf{静态数据集} $\mathcal{D}$，在\textbf{不收集新数据}的情况下，训练一个能最大化奖励的策略 $\pi_\theta$。
\end{importantbox}

Offline RL 的价值：
\begin{itemize}
    \item 利用人类、现有系统已收集的数据。
    \item 在线数据收集可能危险或昂贵。
    \item 复用之前的实验、项目、机器人的数据。
\end{itemize}

\begin{knowledgebox}{Online vs. Offline RL}
Online RL 交替进行"收集数据"和"更新策略"。Offline RL 只有一个静态数据集，训练完成后才部署。也可以先 offline 再 online（混合模式）。
\end{knowledgebox}

\subsection{本章小结}

Offline RL 解决的是"如何在不与环境交互的情况下，从已有数据中学习好的策略"这一问题。

